Dues càrregues elèctriques ($Q_1$ i $Q_2$) estan disposades tal com mostra la figura. Coneixem les dades següents: $Q_1 = 2,00\ \mathrm{\mu C}$, $Q_2 = -4,00\ \mathrm{\mu C}$, $x = 5,00\ \mathrm{m}$ i $d = 3,00\ \mathrm{m}$.

\figura[0.5]{fig1.png}

\begin{apartat}{1}
Representeu i calculeu el camp elèctric (mòdul, direcció i sentit) en el punt $P$, i calculeu també el potencial elèctric en el mateix punt.
\end{apartat}

\begin{apartat}{1}
Canviem les dues càrregues $Q_1$ i $Q_2$ per unes altres amb valors diferents, però situades en la mateixa posició que les originals. Amb aquesta nova configuració, el camp elèctric creat per les dues càrregues sobre el segment $x$ s'anul·la a 1~m de distància de la nova càrrega $Q_1$. Expliqueu raonadament quin serà el signe d'aquestes càrregues i calculeu la relació que hi haurà entre els seus valors.
\end{apartat}

\blocetiquetat{Dada:}{$k = \dfrac{1}{4\pi\varepsilon_0} = 8,99\times 10^{9}\ \mathrm{N\,m^2\,C^{-2}}$.}
